\section{Discussion}

\paragraph{An inversion, not a degradation.}
The central result is not that out-of-distribution detection weakens under adversarial
invariance training but that the ordering reverses. Degradation would be unsurprising
--- a representation stripped of domain information should be worse at recognising a
domain, and chance is what that argument predicts, which is what a single-encoder model
reaches at $\lam = 0.25$ ($0.525$). Going \emph{past} chance is not predicted by it,
and is what Section~\ref{sec:mechanism} accounts for. The results then compose into one
failure sequence: a model trained with invariance to acquisition domain, as
recommended, is deployed at a new site, where it increasingly predicts the benign
majority class --- actinic keratosis recall $0.245 \to 0.016$ --- with higher
confidence than on its own test data, while the monitor that should flag those images
ranks them as more routine than familiar ones and every metric the developer watches
moves favourably.

\paragraph{The failure belongs to gradient reversal, not to invariance.}
Reducing leakage is not sufficient to produce it, and in this sweep leakage is not even
achieved: it bottoms out at $0.555$ against a floor of $0.5$ and rises again at the
largest weights while detection stays inverted. MMD reduces leakage along the same axis
and collapses nothing at the levels it reaches with the classifier intact, where
gradient reversal has already inverted. The two also fail in opposite manners: pushed
harder, MMD becomes less accurate \emph{and} less confident and its out-of-domain
calibration error falls, while under gradient reversal accuracy falls as confidence
rises and calibration error triples. One degrades honestly, the other is confidently
wrong. Wald et al.~\cite{wald2021} showed that enforcing multi-domain calibration
\emph{produces} invariance; the converse fails here, which suggests that where
invariance is the goal, calibration-based or moment-matching objectives may be
preferable because they leave the out-of-domain uncertainty signal intact. We offer
that as a hypothesis: MMD never reached the invariance level gradient reversal reaches
without destroying the task, and we tested one adversarial formulation, so whether CDAN
or adversarial discriminative adaptation behaves the same way is not established here.

\paragraph{The danger zone is where the method succeeds.}
Risk scales with \emph{achieved} invariance, not with the nominal weight or the
presence of an adversarial term. Where the adversary never learned, leakage stayed at
$0.94$ across four orders of magnitude of $\lam$ and the safety failure did not appear
either. A weak adversary produces neither debiasing nor harm --- so many deployed
configurations may be harmless, and a configuration's safety cannot be read off its
hyperparameters.

\paragraph{Tolerance without generalisation.}\label{sec:tolerance}
Full-spectrum OOD detection \cite{yang2023} prescribes that covariate-shifted
in-distribution inputs be accepted rather than flagged. Our models achieve that
prescription and are unsafe because of it: they accept out-of-domain images more
readily than their own test data while classifying them at $0.254$ balanced accuracy
against a floor of $0.167$. Acceptance and competence are separate properties, and
current protocols evaluate them separately or not at all. Any claim of tolerance to
covariate shift should carry cross-domain accuracy on the same inputs, since tolerating
inputs a model cannot classify is a failure mode rather than a success.

\paragraph{Recommendation.}
Report OOD detection and out-of-domain calibration \textbf{as a function of achieved
invariance} rather than at a single operating point; report cross-domain accuracy
alongside any claim of tolerance to covariate shift; and state, when a separate
dataset serves as the OOD set, that this measures covariate shift rather than
semantic novelty. Practitioners tuning $\lam$ already train at several values, so
this changes what is reported rather than what is run. It converts an invisible
failure into a visible number.

\paragraph{Monitor outside the adversarial path.}
Since the collapse is a property of the representation that received adversarial
gradient, one mitigation is to score novelty on a representation that did not. A
frozen ImageNet ResNet-50 gives Mahalanobis AUROC $0.988$ on Fitzpatrick17k and
$0.833$ on \texttt{pad\_heldout}, against the trained model's $0.764$ and $0.843$ at
$\lam = 0$ and $0.391$ and $0.427$ at $\lam = 2$. It never inverts, but \textbf{we
had pre-registered $0.90$ in all four cells and \texttt{pad\_heldout} does not reach
it}, where the frozen encoder merely matches the untrained model rather than
improving on it. A partial mitigation, not a fix; retreating to an earlier layer of
the \emph{same} encoder does not work at all (Section~\ref{sec:nonmediation}).

\paragraph{The inversion is deepened, not created, in the minority classes.}
\label{sec:perclass}
Per-class AUROC, which unlike the pooled value does not depend on the ID class mix,
is $0.14$--$0.39$ on \texttt{pad\_heldout} and $0.10$--$0.34$ on Fitzpatrick17k at
$\lam = 2$ for every non-nevus class. Two are already at or below chance beforehand: on
\texttt{pad\_heldout} at $\lam = 0$, actinic keratosis is $0.499$ and squamous cell
carcinoma $0.444$.
Gradient reversal deepens an inversion that already exists for the rarest and most
clinically consequential classes, and the pooled AUROC at $\lam = 0$ conceals this
because nevus is half the in-distribution set.

\paragraph{The mechanism, and what remains open.}
Gradient reversal removes the acquisition signal by moving out-of-domain images into
the region the majority class occupies (Section~\ref{sec:mechanism}) --- an efficient
way to satisfy the adversary, since the densest part of the latent space is where
domain information is hardest to recover, and simultaneously the worst outcome for a
distance-based novelty score, which reads density as familiarity. The two objectives
are not merely in tension: invariance is obtained by the move that defeats detection.
What remains open is why gradient reversal takes this route and moment matching does
not, why the synthetic reproduction failed, and whether the majority class is special
or merely the densest available --- a dataset with two comparably sized dominant
classes would separate the last of these.

\section{Limitations}
\label{sec:limitations}

\textbf{Absolute cross-domain performance is poor for every model here.} Plain
accuracy on \texttt{pad\_heldout} is $0.245$ for the factorized model and $0.391$ for
the baseline against a $0.385$ majority rate, so no model here is deployable. Our
claim concerns the \emph{effect of $\lam$}, and the deployment pathway above is a
mechanism of harm rather than a description of a fielded system.

\textbf{The inversion is observed only across imaging modality.} ISIC is dermoscopy
while PAD-UFES and Fitzpatrick17k are clinical photography, so the main contrast
confounds site with modality and Fitzpatrick17k does not resolve it. The
same-modality contrast within ISIC (Section~\ref{sec:site}) separates the claims: the
collapse reproduces between two hospitals, the inversion does not. We cannot say
whether that reflects the size of the shift, modality gaps specifically, or these
datasets.

Three attempts to replicate outside dermatology returned no result, for the same
underlying reason in different guises: the manipulation did not take, so there was
nothing to measure. On three-hospital Camelyon17 the adversary head never trained; on
iWildCam it trained but the encoder never became invariant; and on Camelyon17
restricted to \emph{two} hospitals --- the obvious diagnosis, matching the dermatology
domain count, with the gradient path verified first --- the adversary reached chance
($\ln 2$) while a post-hoc probe still read the hospital at $0.92$--$0.95$, a fall of
at most $0.037$ against the $0.10$ we had pre-committed. The pre-registered stopping
rule returns a negative and we stop. \textbf{Domain count is not the barrier}, and we
cannot say what is; the pattern --- an adversary driven to chance while the information
it was asked to remove stays linearly decodable --- is a known failure mode of
adversarial removal \cite{elazar2018}. \textbf{We report no inversion outside
dermatology}; a single-seed sub-chance reading there did not replicate
($0.534 \pm 0.050$ over three seeds). What those runs do support is the dissociation in
Section~\ref{sec:nonmediation}. Details in Appendix~\ref{app:external}.

Generality beyond dermatology remains untested. \textbf{Backbone coverage.} Comparable invariance was not reached with a ResNet-50
backbone, so whether the inversion extends beyond EfficientNet-B3 is untested. Three of
the nine runs failed their adversary gate outright and are reported as null experiments
rather than negative results. \textbf{The objective comparison is bounded.} MMD never reached the invariance
level gradient reversal attains; it could only approach it by destroying the
classifier. We therefore cannot say what moment-matching would do at matched
invariance, only that at the leakage levels it reaches, gradient reversal has
already inverted and MMD has not.

\textbf{Detector non-independence.} The five scoring rules are functions of one
16-dimensional latent; the backbone analysis and the single-encoder replication carry
the independent evidence. \textbf{Calibration arithmetic.} Out-of-domain ECE is close
to mean confidence minus accuracy here, so those three results are one phenomenon
viewed three ways. \textbf{Seeds.} Five at every $\lam$ on the leak-free split, and
three at two interior points of the image-level sweep in the appendix; effect sizes
are large relative to seed variance, but we claim no formal test across seven
points.

\textbf{Splits.} The main sweep uses the image-level ISIC split, in which
$3{,}041$ of $5{,}067$ test images ($60.0\%$) share a \texttt{lesion\_id} with
training or validation. This inflates in-distribution balanced accuracy by
$0.20$ ($0.496 \to 0.692$ at $\lam = 0$), so the absolute values in
Table~\ref{tab:dial} are optimistic by that amount;
Section~\ref{sec:lesionsplit} reports the lesion-disjoint replication at five
seeds, under which every conclusion is unchanged. PAD splits are patient-level,
disjoint by patient and by lesion. \textbf{Third-domain coverage.} 3{,}887 of 16{,}577 Fitzpatrick17k images were
obtained; the set is web-scraped under mixed copyright, used for evaluation only,
and not redistributed.

% ------------------------------------------------------------
\section{Conclusion}

Under domain-adversarial training, out-of-distribution detection does not merely
weaken. It reverses: a model asked which inputs are unfamiliar answers with its
own test data and ranks a new clinic's images as the more routine of the two.
The reversal appears at the smallest adversarial weight that achieves
invariance, holds on a dataset the adversary never saw, reproduces without the
factorized architecture, and does not appear under moment-matching objectives
that reduce leakage along the same axis.

It is invisible to everything a developer monitors. In-distribution accuracy
holds, in-distribution calibration is unchanged and the domain leakage probe
improves exactly as intended --- under an image-level and a lesion-disjoint ISIC
split alike --- while confidence on unseen domains overtakes confidence on
familiar data and recall for a premalignant class falls to $0.016$.

The cause is identifiable: invariance is obtained by moving out-of-domain
images into the majority class's region of the latent space, which is exactly
the move that makes a density-based novelty score call them familiar. Acting on
it does not require the explanation. Reporting out-of-distribution detection,
out-of-domain calibration and cross-domain accuracy as functions of achieved
invariance costs almost nothing and makes the failure visible.

% ------------------------------------------------------------
